\section{Requisitos del Sistema}

%Los requisitos del sistema se derivan de las dificultades de estandarización, descripción y recuperación de la documentación histórica institucional. Su formulación corresponde a una etapa previa al diseño de la solución, por lo que establece las capacidades esperadas y las condiciones generales de operación sin definir componentes, tecnologías o mecanismos de implementación.

\subsection{Requisitos Funcionales}

\begin{enumerate}

    \item[\textbf{RF-01}] El sistema deberá permitir obtener el contenido textual
    presente en documentos digitalizados.

    \item[\textbf{RF-02}] El sistema deberá conservar, cuando la naturaleza del
    documento lo permita, el orden de lectura y la organización lógica de elementos
    como títulos, párrafos, columnas y tablas.

    \item[\textbf{RF-03}] El sistema deberá transformar la información recuperada
    en una representación procesable que permita identificar el contenido de cada
    documento.

    \item[\textbf{RF-04}] El sistema deberá generar información complementaria,
    como metadatos y descripciones de contenido, que apoye la homogeneización,
    indexación y búsqueda de los documentos.

    \item[\textbf{RF-05}] El sistema deberá admitir el análisis de documentos con
    diferentes estructuras, condiciones de digitalización y niveles de degradación
    visual.

    \item[\textbf{RF-06}] El proyecto deberá permitir comparar distintas
    alternativas de procesamiento documental mediante documentos representativos
    y criterios comunes de evaluación.

\end{enumerate}

\subsection{Requisitos No Funcionales}

\begin{enumerate}

    \item[\textbf{RNF-01}] La información generada deberá mantener fidelidad
    respecto del documento original, minimizando sustituciones, omisiones,
    incorporaciones de contenido y alteraciones estructurales.

    \item[\textbf{RNF-02}] La solución deberá utilizar eficientemente la capacidad
    de procesamiento y memoria disponible en la infraestructura institucional.

    \item[\textbf{RNF-03}] El sistema deberá permitir el procesamiento de un
    volumen elevado de documentos sin comprometer su estabilidad ni la continuidad
    operacional.

    \item[\textbf{RNF-04}] La implementación deberá priorizar herramientas libres
    y de código abierto que puedan ejecutarse dentro de la infraestructura
    institucional.

    \item[\textbf{RNF-05}] La selección de la alternativa de procesamiento deberá
    considerar conjuntamente la fidelidad de los resultados, el tiempo de
    procesamiento, el consumo de recursos y la estabilidad de ejecución.

    \item[\textbf{RNF-06}] Las comparaciones experimentales deberán ser
    reproducibles y mantener condiciones equivalentes para los documentos,
    configuraciones y recursos evaluados.

\end{enumerate}

\subsection{Requisitos Ambientales}

% Los requisitos ambientales describen las condiciones generales que deberá satisfacer el entorno de ejecución. Las capacidades mínimas definitivas deberán determinarse posteriormente mediante la evaluación experimental de las alternativas consideradas.

\subsubsection{Hardware}

\begin{enumerate}

    \item[\textbf{RA-H01}] La solución deberá poder ejecutarse utilizando la
    infraestructura computacional disponible en la institución, caracterizada por
    una capacidad limitada de procesamiento y memoria.

    \item[\textbf{RA-H02}] El sistema no deberá depender obligatoriamente de
    hardware especializado que no se encuentre disponible en el entorno
    institucional.

    \item[\textbf{RA-H03}] El entorno deberá disponer de capacidad suficiente para
    almacenar los documentos de entrada, los modelos de procesamiento y los
    resultados generados.

\end{enumerate}

\subsubsection{Software}

\begin{enumerate}

    \item[\textbf{RA-S01}] El entorno deberá admitir herramientas de procesamiento
    de imágenes, reconocimiento textual y análisis de estructura documental.

    \item[\textbf{RA-S02}] Las herramientas seleccionadas deberán ser compatibles
    con el sistema operativo y el hardware disponible en la institución.

    \item[\textbf{RA-S03}] El entorno deberá permitir la ejecución local de los
    modelos y mecanismos de procesamiento seleccionados.

    \item[\textbf{RA-S04}] El entorno de evaluación deberá permitir registrar y
    comparar la calidad de los resultados, los tiempos de ejecución y el consumo
    de recursos de cada alternativa.

    \item[\textbf{RA-S05}] Se deben ocupar herramientas de desarrollo utilizadas por el equipo de desarrollo de la institución.

\end{enumerate}